\documentclass[11pt,a4paper]{article}

% ---------- Packages ----------
\usepackage[margin=1in]{geometry}
\usepackage{amsmath, amssymb}
\usepackage{booktabs}
\usepackage{array}
\usepackage{siunitx}
\usepackage{graphicx}
\usepackage{float}
\usepackage{hyperref}
\usepackage{xcolor}

% ---------- Document ----------
\begin{document}

% ============================================================
% Title Page
% ============================================================

\begin{titlepage}
    \centering

    \vspace*{2cm}

    {\Large \textbf{ELL710 Coding Theory}\par}

    \vspace{1.5cm}

    {\huge \textbf{A Design Study}\par}

    \vspace{0.5cm}

    {\LARGE Telemetry Downlink for a CubeSat in Low Earth Orbit\par}

    \vspace{2cm}

    % \begin{tabular}{ll}
    %     \textbf{Course:} & ELL710 / ELL7406 Coding Theory \\
    %     \textbf{Department:} & Department of Electrical Engineering \\
    %     \textbf{Institute:} & Indian Institute of Technology Delhi \\
    %     \textbf{Module:} & Module 1 -- Source Coding \\
    %     \textbf{Due Date:} & 31 August 2026
    % \end{tabular}

    \vfill

\end{titlepage}


% ============================================================
% Introduction / Problem Setup
% ============================================================

\section{Problem Setup}

We consider a 3U CubeSat operating in a low Earth orbit (LEO)
that continuously generates housekeeping telemetry. Each telemetry
record belongs to one of eight possible event types. The records
accumulate in an onboard buffer and are downlinked to a single ground
station during the one visible pass available to the spacecraft in
each orbit.

The communication link has a modest data rate and a short visibility
window. Any data that cannot be transmitted during a pass is lost,
since new telemetry continues to arrive during the following orbit.
Therefore, efficient source coding is required to reduce the number of
bits needed to transmit the telemetry.

\subsection{System Parameters}

The relevant system parameters are summarized in Table~\ref{tab:parameters}.

\begin{table}[H]
    \centering
    \caption{CubeSat telemetry downlink parameters.}
    \label{tab:parameters}
    \begin{tabular}{lc}
        \toprule
        \textbf{Parameter} & \textbf{Value} \\
        \midrule
        Downlink & UHF \\
        Data rate & \SI{9600}{bit/s} \\
        Usable pass duration & \SI{480}{s} \\
        Records buffered between passes & 1,700,000 \\
        Sunlit fraction of orbit & 65\% \\
        Eclipse fraction of orbit & 35\% \\
        \bottomrule
    \end{tabular}
\end{table}

The capacity of one downlink pass is therefore

\begin{equation}
    C_{\mathrm{pass}}
    = 9600 \times 480
    = 4{,}608{,}000 \text{ bits}.
\end{equation}

Thus, the source coding scheme must be sufficiently efficient to
transmit the buffered telemetry within this available capacity.

\subsection{Telemetry Event Distributions}

The telemetry source consists of eight event types. Their probabilities
depend on the operating mode of the spacecraft. The source is assumed
to be memoryless within each mode, and the operating mode remains
constant within a frame.

\begin{table}[H]
    \centering
    \caption{Telemetry event probabilities in the two operating modes.}
    \label{tab:probabilities}
    \begin{tabular}{lcc}
        \toprule
        \textbf{Event} & \textbf{Sunlit} & \textbf{Eclipse} \\
        \midrule
        E1 -- Housekeeping OK       & 0.36 & 0.18 \\
        E2 -- Attitude update       & 0.20 & 0.10 \\
        E3 -- Thermal reading       & 0.14 & 0.24 \\
        E4 -- Power and battery     & 0.11 & 0.26 \\
        E5 -- Payload status        & 0.08 & 0.06 \\
        E6 -- Link event            & 0.05 & 0.08 \\
        E7 -- Fault warning         & 0.04 & 0.05 \\
        E8 -- Critical fault        & 0.02 & 0.03 \\
        \bottomrule
    \end{tabular}
\end{table}

The spacecraft operates in sunlight for 65\% of the orbit and in
eclipse for 35\%. When a single probability distribution is required,
the mode-averaged distribution is defined as

\begin{equation}
    p_m = 0.65p_s + 0.35p_e,
\end{equation}

where \(p_s\) and \(p_e\) denote the sunlit and eclipse distributions,
respectively.

\subsection{Coding Objective}

The objective of this study is to investigate different source coding
schemes for the telemetry and determine which design is most suitable
for the CubeSat downlink.

The analysis considers binary Huffman coding, the overhead associated
with transmitting a code table, separate codes for different operating
modes, coding with a mismatched source model, ternary Huffman coding,
and Shannon--Fano--Elias coding. These designs are ultimately compared
in terms of the total number of bits required to transmit the buffered
telemetry and their utilization of the available downlink capacity.

\newpage

% ============================================================
% Part A -- Get a Baseline
% ============================================================

\section{Part A -- Get a Baseline}

A binary Huffman code is constructed using the sunlit probability
distribution. The two least-probable symbols are repeatedly merged
until a single tree is obtained.

\subsection{Huffman Tree}

The resulting Huffman tree is shown in Figure~\ref{fig:huffman_sunlit}.
The left and right branches correspond to binary symbols 0 and 1,
respectively.

\begin{figure}[H]
    \centering
    % Insert your Huffman tree here
    \includegraphics[width=0.65\textwidth]{huffman_sunlit.png}
    \caption{Binary Huffman tree for the sunlit distribution.}
    \label{fig:huffman_sunlit}
\end{figure}

\subsection{Codewords and Lengths}

The resulting codewords and their lengths are shown in
Table~\ref{tab:huffman_codes}.

\begin{table}[H]
    \centering
    \caption{Sunlit Huffman code.}
    \label{tab:huffman_codes}
    \begin{tabular}{lcc}
        \toprule
        \textbf{Event} & \textbf{Codeword} & \textbf{Length} \\
        \midrule
        E1 & 11 & 2 \\
        E2 & 01 & 2 \\
        E3 & 101 & 3 \\
        E4 & 100 & 3 \\
        E5 & 000 & 3 \\
        E6 & 0010 & 4 \\
        E7 & 00111 & 5 \\
        E8 & 00110 & 5 \\
        \bottomrule
    \end{tabular}
\end{table}

Thus, the length vector is

\begin{equation}
    \mathbf{l} = (l_1,l_2,\ldots,l_8)
    = (2, 2, 3, 3, 3, 4, 5, 5).
\end{equation}

\subsection{Expected Code Length}

The expected number of bits per telemetry record is

\begin{equation}
    L = \sum_{i=1}^{8} p_i l_i
    = \boxed{2.6100\ \text{bits/record}}.
\end{equation}

\subsection{Entropy and Efficiency}

The entropy of the sunlit distribution is

\begin{equation}
    H = -\sum_{i=1}^{8}p_i\log_2 p_i
    = \boxed{2.55\ \text{bits/record}}.
\end{equation}

The coding efficiency is therefore

\begin{equation}
    \eta = \frac{H}{L}\times100
    = \boxed{97.65\%}.
\end{equation}

\subsection{Kraft Inequality}

For a binary prefix code, the codeword lengths must satisfy

\begin{equation}
    \sum_{i=1}^{8}2^{-l_i}\leq1.
\end{equation}

For the obtained code,

\begin{equation}
    \sum_{i=1}^{8}2^{-l_i}
    = 1.000 \leq 1,
\end{equation}

so the Kraft inequality is satisfied.

\subsection{Comparison with Fixed-Length Coding}

A fixed-length representation requires

\[
\log_2 8 = 3
\]

bits per record. The Huffman code therefore saves

\begin{equation}
    3 - 2.6100 = \boxed{0.3900\ \text{bits/record}}
\end{equation}

on average compared with the fixed-length code.

% ============================================================
% Part B -- Pay for the Code Table
% ============================================================

\section{Part B -- Pay for the Code Table}

Each frame carries its own canonical Huffman code table. Since there
are eight codeword lengths and each length is represented using four
bits, the header requires

\begin{equation}
    H_{\mathrm{header}} = 8 \times 4 = 32\text{ bits}.
\end{equation}

The fixed-length code requires no header since its mapping is fixed in
advance.

\subsection{Expected Bits per Frame}

For a frame containing \(N\) records, the fixed-length code requires

\begin{equation}
    B_{\mathrm{fixed}}(N) = 3N.
\end{equation}

Using the expected Huffman length \(L=2.6100\) bits/record, the Huffman
code requires

\begin{equation}
    B_{\mathrm{Huffman}}(N)
    = 32 + 2.6100N.
\end{equation}

\subsection{Break-even Frame Size}

Huffman coding becomes cheaper when

\begin{align}
    32 + 2.6100N &< 3N,\\
    32 &< 0.3900N.
\end{align}

Thus,

\begin{equation}
    N > \frac{32}{0.3900}
    = 82.05.
\end{equation}

Therefore, the smallest whole-number frame size for which Huffman
coding is cheaper is

\begin{equation}
    \boxed{N=83}.
\end{equation}

At \(N=83\),

\[
B_{\mathrm{fixed}}=249.00\text{ bits},
\qquad
B_{\mathrm{Huffman}}=248.63\text{ bits}.
\]

Hence, Huffman coding is cheaper by \(0.37\) bits per frame at the
break-even point.

\subsection{Saving per Record}

The net saving per record for a frame of size \(N\) is

\begin{equation}
    S(N)
    = \frac{B_{\mathrm{fixed}}(N)-B_{\mathrm{Huffman}}(N)}{N}.
\end{equation}

Substituting the expressions above,

\begin{equation}
    S(N)
    = 0.3900-\frac{32}{N}.
\end{equation}

The savings for the specified frame sizes are shown in
Table~\ref{tab:frame_savings}.

\begin{table}[H]
    \centering
    \caption{Comparison of fixed-length and Huffman coding for different frame sizes.}
    \label{tab:frame_savings}
    \begin{tabular}{cccc}
        \toprule
        \(N\) & Fixed-length (bits) & Huffman (bits) &
        Saving (bits/record) \\
        \midrule
        50  & 150.0  & 162.5  & -0.250 \\
        200 & 600.0  & 554.0  & 0.230 \\
        500 & 1500.0 & 1337.0 & 0.326 \\
        \bottomrule
    \end{tabular}
\end{table}

For \(N=50\), the header overhead exceeds the compression benefit, so
Huffman coding actually requires more bits per record. At \(N=200\)
and \(N=500\), Huffman coding provides a net saving.

As \(N\) increases,

\begin{equation}
    \lim_{N\rightarrow\infty} S(N)
    = 0.3900\text{ bits/record}.
\end{equation}

Thus, the header overhead becomes negligible for sufficiently large
frames. A frame size in the range of 200--500 records is therefore a
safe choice based on the coding overhead considered here.

\subsection{Maximum Huffman Codeword Length}

For a binary Huffman tree with eight symbols, the most unbalanced
possible tree can have a leaf at depth \textbf{7}.

\begin{figure}[H]
    \centering
    % Insert maximum-depth Huffman tree here
    \includegraphics[width=0.75\textwidth]{max_depth_tree.png}
    \caption{A maximally unbalanced binary Huffman tree with eight symbols.}
    \label{fig:max_depth_tree}
\end{figure}

Hence, the maximum possible codeword length is 7 bits. Since the
length can range up to 7, a 3-bit field is sufficient to represent
the codeword length. Therefore, the 4-bit field used in the assignment
has one redundant bit per codeword length.

% ============================================================
% Part C -- One Code, or Two?
% ============================================================

\section{Part C -- One Code, or Two?}

Two coding schemes are compared: a single Huffman code constructed
from the mode-averaged distribution, and separate Huffman codes for
the sunlit and eclipse modes. The latter requires only a one-bit mode
flag per frame since both codebooks are known to the transmitter and
receiver in advance.

\subsection{Mode-Averaged Distribution}

The mode-averaged distribution is calculated as

\begin{equation}
    p_m = 0.65p_s + 0.35p_e.
\end{equation}

This gives

\begin{table}[H]
    \centering
    \caption{Mode-averaged telemetry distribution.}
    \label{tab:mode_avg}
    \begin{tabular}{lc}
        \toprule
        \textbf{Event} & \(\boldsymbol{p_m}\) \\
        \midrule
        E1 & 0.2970 \\
        E2 & 0.1650 \\
        E3 & 0.1750 \\
        E4 & 0.1625 \\
        E5 & 0.0730 \\
        E6 & 0.0605 \\
        E7 & 0.0435 \\
        E8 & 0.0235 \\
        \bottomrule
    \end{tabular}
\end{table}

\subsection{Expected Lengths}

A Huffman code constructed from the mode-averaged distribution has an
expected length of

\begin{equation}
    L_{\mathrm{single}} = 2.7225
    \text{ bits/record}.
\end{equation}

For the two-code design, the expected lengths within each mode are

\begin{align}
    L_{\mathrm{sunlit}} &= 2.6100
    \text{ bits/record},\\
    L_{\mathrm{eclipse}} &= 2.7200
    \text{ bits/record}.
\end{align}

Averaging these according to the fraction of time spent in each mode,

\begin{align}
    L_{\mathrm{two}}
    &= 0.65L_{\mathrm{sunlit}}
       + 0.35L_{\mathrm{eclipse}}\\
    &= 0.65(2.6100)+0.35(2.7200)\\
    &= \boxed{2.6485\text{ bits/record}}.
\end{align}

Thus, using separate Huffman codes reduces the expected coding length
compared with the single-code design.

\subsection{Saving of the Two-Code Design}

The saving of the two-code design over the single-code design is

\begin{equation}
    \Delta L
    = L_{\mathrm{single}}-L_{\mathrm{two}}
    = 2.7225-2.6485
    = \boxed{0.0740\text{ bits/record}}.
\end{equation}

The percentage saving is

\begin{equation}
    \frac{\Delta L}{L_{\mathrm{single}}}\times100
    = \boxed{2.7181\%}.
\end{equation}

\subsection{Minimum Frame Size}

For a frame containing \(N\) records, the single-code design requires

\begin{equation}
    B_{\mathrm{single}} = NL_{\mathrm{single}},
\end{equation}

while the two-code design requires one additional mode-flag bit:

\begin{equation}
    B_{\mathrm{two}} = NL_{\mathrm{two}}+1.
\end{equation}

The two-code design becomes cheaper when

\begin{equation}
    NL_{\mathrm{two}}+1 < NL_{\mathrm{single}}.
\end{equation}

Using the calculated expected lengths,

\begin{equation}
    N > \frac{1}{L_{\mathrm{single}}-L_{\mathrm{two}}}
    = \frac{1}{0.0740}
    \approx 13.51.
\end{equation}

Therefore, the smallest whole-number frame size for which the
two-code design is cheaper is

\begin{equation}
    \boxed{N=14}.
\end{equation}

% ============================================================
% Part D -- Put a Number on Using the Wrong Model
% ============================================================

\section{Part D -- Put a Number on Using the Wrong Model}

We analyze the performance of the single-code design (optimized for the mode-averaged distribution \(p_m\)) when the spacecraft is in eclipse (operating under the eclipse distribution \(p_e\)).

The expected length achieved by this mismatched code is:
\begin{equation}
    L_{\mathrm{wrong}} = \sum_{i=1}^{8} p_e(E_i) l_m(E_i) = \boxed{2.8200\ \text{bits/record}}.
\end{equation}

This length is compared with two reference values:
\begin{itemize}
    \item \textbf{Eclipse-optimal Huffman code}: \(L_{\mathrm{eclipse}} = 2.7200\) bits/record, yielding a penalty of:
    \begin{equation}
        \Delta L_{\mathrm{optimal}} = L_{\mathrm{wrong}} - L_{\mathrm{eclipse}} = \boxed{0.1000\ \text{bits/record}}.
    \end{equation}
    \item \textbf{Eclipse entropy}: \(H(p_e) \approx 2.6798\) bits/record, yielding a penalty of:
    \begin{equation}
        \Delta L_{\mathrm{entropy}} = L_{\mathrm{wrong}} - H(p_e) \approx \boxed{0.1402\ \text{bits/record}}.
    \end{equation}
\end{itemize}

The Kullback--Leibler (KL) divergence measuring the theoretical mismatch penalty is:
\begin{equation}
    D(p_e \parallel p_m) = \sum_{i=1}^{8} p_e(E_i) \log_2 \left(\frac{p_e(E_i)}{p_m(E_i)}\right) \approx \boxed{0.1193\ \text{bits/record}}.
\end{equation}

Comparing the two, the penalty relative to entropy is larger than the KL divergence (\(0.1402 > 0.1193\)). This difference arises because the KL divergence assumes ideal fractional codeword lengths (\(l_m^*(x) = -\log_2 p_m(x)\)), whereas a real Huffman code must use integer lengths. The exact relationship is:
\begin{equation}
    L_{\mathrm{wrong}} - H(p_e) = D(p_e \parallel p_m) + \sum_{x} p_e(x) r(x),
\end{equation}
where \(r(x) = l_m(x) + \log_2 p_m(x)\) is the pointwise Huffman redundancy. The difference consists of the expected discretization overhead under the true distribution \(p_e\):
\begin{equation}
    \text{Expected Discretization Overhead under } p_e = \sum_{x} p_e(x) r(x) \approx \boxed{0.0209\ \text{bits/record}}.
\end{equation}
For reference, the discretization overhead (redundancy) of this single code under its own design distribution is \(L_{\mathrm{single}} - H(p_m) \approx 0.0630\) bits/record.

% ============================================================
% Part E -- Consider a Ternary Link
% ============================================================

\section{Part E -- Consider a Ternary Link}

A proposed radio upgrade would use a three-level line code, making the channel alphabet size \(D = 3\). We construct a ternary Huffman code for the mode-averaged distribution \(p_m\).

\subsection{Dummy Symbols and Divisibility Condition}

In a \(D\)-ary Huffman tree, every merge combines \(D\) nodes into a single node, which reduces the number of active nodes in the pool by \(D - 1\). To reach exactly one root node starting from \(N_s\) source symbols and \(N_d\) dummy symbols, the total number of initial nodes must satisfy the divisibility condition:
\begin{equation}
    (N_s + N_d - 1) \equiv 0 \pmod{D - 1}.
\end{equation}

For \(N_s = 8\) and \(D = 3\):
\begin{equation}
    (8 + N_d - 1) \equiv (7 + N_d) \equiv 0 \pmod{2}.
\end{equation}

The smallest non-negative integer satisfying this condition is \(\boxed{N_d = 1}\) dummy symbol (with probability \(0\)).

\subsection{Ternary Huffman Code Properties}

Using the \(D\)-ary Huffman encoder, the resulting ternary codewords and their lengths are summarized in Table~\ref{tab:ternary_codes}.

\begin{table}[H]
    \centering
    \caption{Ternary Huffman code for the mode-averaged distribution \(p_m\).}
    \label{tab:ternary_codes}
    \begin{tabular}{lcc}
        \toprule
        \textbf{Event} & \textbf{Codeword} & \textbf{Length (symbols)} \\
        \midrule
        E1 & \texttt{1}   & 1 \\
        E2 & \texttt{21}  & 2 \\
        E3 & \texttt{22}  & 2 \\
        E4 & \texttt{20}  & 2 \\
        E5 & \texttt{02}  & 2 \\
        E6 & \texttt{00}  & 2 \\
        E7 & \texttt{012} & 3 \\
        E8 & \texttt{011} & 3 \\
        \bottomrule
    \end{tabular}
\end{table}

This yields the length vector:
\begin{equation}
    \mathbf{l} = (1, 2, 2, 2, 2, 2, 3, 3).
\end{equation}

The performance metrics of this ternary code are:
\begin{itemize}
    \item \textbf{Expected ternary length \(L_{\mathrm{ternary}}\)}:
    \begin{equation}
        L_{\mathrm{ternary}} = \sum_{i=1}^{8} p_m(E_i) l_i = \boxed{1.7700\ \text{ternary symbols/record}}.
    \end{equation}
    \item \textbf{Ternary entropy \(H_3(p_m)\)}:
    \begin{equation}
        H_3(p_m) = -\sum_{i=1}^{8} p_m(E_i) \log_3 p_m(E_i) \approx \boxed{1.6780\ \text{ternary symbols/record}}.
    \end{equation}
    \item \textbf{Coding efficiency \(\eta_3\)}:
    \begin{equation}
        \eta_3 = \frac{H_3(p_m)}{L_{\mathrm{ternary}}} \approx \boxed{94.80\%}.
    \end{equation}
\end{itemize}

\subsection{Comparison: Ternary vs. Binary Link}

To determine if the ternary link is superior to the binary one, they must be compared on a common physical basis. The conclusion depends on the physical limits of the channel:

\begin{enumerate}
    \item \textbf{Baud-rate (symbol-rate) limited channel}: 
    If the transmitter is bandwidth-limited (i.e., can transmit a fixed number of symbols per second \(R_s\)), the ternary link is \textbf{better}. The binary code requires \(L_{\mathrm{binary}} = 2.7225\) bits/record (throughput of \(0.3673\) records/symbol), whereas the ternary code requires \(L_{\mathrm{ternary}} = 1.7700\) symbols/record (throughput of \(0.5650\) records/symbol). This represents a throughput increase of \(\boxed{53.81\%}\).
    
    \item \textbf{Power-limited channel}: 
    If the channel is limited by transmitter power and noise, the binary link is \textbf{better}. Transitioning from a 2-level line code (e.g., BPSK) to a 3-level line code (e.g., 3-PAM) requires a higher signal-to-noise ratio (SNR) to maintain the same bit error rate. Under fixed power and bandwidth constraints, the Shannon capacity (in bits/sec) is constant. Since the binary Huffman code is already highly efficient (\(97.65\%\)), there is negligible coding gain to be made by migrating to ternary signaling, and the binary link remains preferable due to its simpler hardware and lower SNR requirement.
\end{enumerate}
% ============================================================
% Part F -- Try Shannon-Fano-Elias
% ============================================================

\section{Part F -- Try Shannon--Fano--Elias}

We construct the Shannon--Fano--Elias (SFE) code for the mode-averaged distribution $p_m$. The codeword length for each symbol $x$ is determined by:
\begin{equation}
    l(x) = \left\lceil \log_2 \frac{1}{p_m(x)} \right\rceil + 1.
\end{equation}
The modified cumulative distribution function $\bar{F}(x)$ is computed as:
\begin{equation}
    \bar{F}(x_i) = \sum_{j=1}^{i-1} p_m(x_j) + \frac{1}{2} p_m(x_i).
\end{equation}
The binary codeword for each event is obtained by taking the first $l(x)$ bits of the fractional binary expansion of $\bar{F}(x)$.

\subsection{SFE Code Construction Table}

The cumulative probabilities, modified CDF values, codeword lengths, and binary codewords are summarized in Table~\ref{tab:sfe_code}.

\begin{table}[H]
    \centering
    \caption{Shannon--Fano--Elias code table for the mode-averaged distribution $p_m$.}
    \label{tab:sfe_code}
    \begin{tabular}{lccccc}
        \toprule
        \textbf{Event} & $\boldsymbol{p_m(x)}$ & $\boldsymbol{F(x)}$ & $\boldsymbol{\bar{F}(x)}$ & $\boldsymbol{l(x)}$ & \textbf{Codeword} \\
        \midrule
        E1 & 0.2970 & 0.2970 & 0.148500 & 3 & \texttt{001} \\
        E2 & 0.1650 & 0.4620 & 0.379500 & 4 & \texttt{0110} \\
        E3 & 0.1750 & 0.6370 & 0.549500 & 4 & \texttt{1000} \\
        E4 & 0.1625 & 0.7995 & 0.718250 & 4 & \texttt{1011} \\
        E5 & 0.0730 & 0.8725 & 0.836000 & 5 & \texttt{11010} \\
        E6 & 0.0605 & 0.9330 & 0.902750 & 6 & \texttt{111001} \\
        E7 & 0.0435 & 0.9765 & 0.954750 & 6 & \texttt{111101} \\
        E8 & 0.0235 & 1.0000 & 0.988250 & 7 & \texttt{1111110} \\
        \bottomrule
    \end{tabular}
\end{table}

\subsection{Performance Metrics and Bound Verification}

The expected codeword length achieved by the SFE code is:
\begin{equation}
    L_{\mathrm{SFE}} = \sum_{i=1}^{8} p_m(E_i) l(E_i) = \boxed{4.0545\ \text{bits/record}}.
\end{equation}

The entropy of the mode-averaged distribution is $H(p_m) \approx 2.6595\text{ bits/record}$. 

We verify whether the expected length satisfies the theoretical SFE bound $H(p_m) + 1 \le L_{\mathrm{SFE}} < H(p_m) + 2$:
\begin{equation}
    3.6595 \le 4.0545 < 4.6595.
\end{equation}
The calculated expected length of $4.0545$ bits/record strictly satisfies the theoretical bounds.

\subsection{Downlink Capacity Analysis}

To downlink all $1,700,000$ buffered records using the SFE design, the total required transmission capacity is:
\begin{equation}
    B_{\mathrm{SFE}} = 1,700,000 \times 4.0545 = \boxed{6,892,650\ \text{bits}}.
\end{equation}

Given the available single-pass capacity of $C_{\mathrm{pass}} = 4,608,000\text{ bits}$, the SFE code requires:
\begin{equation}
    \frac{6,892,650}{4,608,000} \times 100 \approx \boxed{149.58\%}
\end{equation}
of the downlink capacity.

\textbf{Candidate Assessment:}
Shannon--Fano--Elias coding is \textbf{not a serious candidate} for this downlink system. Because it requires nearly $150\%$ of the available pass capacity, the satellite cannot clear its buffered telemetry during a single pass, resulting in buffer overflow and catastrophic data loss for subsequent orbits. The $+1$ bit addition per symbol required to guarantee prefix-free status introduces significant redundancy ($1.3950$ bits/record above entropy), making SFE primarily a theoretical construction used to introduce arithmetic coding rather than a practical block code.
% ============================================================
% Part G -- Decide What to Fly
% ============================================================

\section{Part G -- Decide What to Fly}

We compare the four candidate source coding designs for downlinking the $1,700,000$ buffered telemetry records over the single visible ground station pass ($C_{\mathrm{pass}} = 4,608,000\text{ bits}$). For framed schemes, an operational frame size of $N = 200$ records is assumed.

\subsection{Quantitative Comparison of Coding Schemes}

The total bits required to transmit the entire telemetry buffer and the corresponding pass capacity utilization for each design are summarized in Table~\ref{tab:final_comparison}.

\begin{table}[H]
    \centering
    \caption{Downlink volume and capacity utilization across candidate designs.}
    \label{tab:final_comparison}
    \begin{tabular}{lcccc}
        \toprule
        \textbf{Design Scheme} & \textbf{Expected Length $L$} & \textbf{Total Bits} & \textbf{Capacity Used} & \textbf{Feasible?} \\
        & (bits/record) & & (\%) & \\
        \midrule
        Fixed-length (3-bit) & 3.0000 & 5,100,000 & 110.68\% & No (Overflow) \\
        Single Huffman ($p_m$) & 2.7225 & 4,628,250 & 100.44\% & No (Overflow) \\
        \textbf{Two-Code Huffman} & \textbf{2.6485} & \textbf{4,510,950} & \textbf{97.89\%} & \textbf{Yes} \\
        Shannon--Fano--Elias (SFE) & 4.0545 & 6,892,650 & 149.58\% & No (Overflow) \\
        \bottomrule
    \end{tabular}
\end{table}

\textit{Note on Two-Code Huffman calculation:} Incorporating a 1-bit mode flag per frame ($N = 200$) requires $8,500$ frame header bits across $8,500$ frames. The overall expected bit rate is $2.6485 + \frac{1}{200} = 2.6535\text{ bits/record}$, yielding a total volume of $4,510,950\text{ bits}$.

\subsection{Design Selection and Justification}

I recommend flying the \textbf{Two-Code Huffman design with a 1-bit frame header flag}.

As shown in Table~\ref{tab:final_comparison}, it is the \textbf{only design that fits within the downlink capacity constraint} of $4.608\text{ Mbit}$. The fixed-length code ($110.68\%$), single Huffman code ($100.44\%$), and SFE code ($149.58\%$) all exceed $100\%$ capacity, meaning telemetry would accumulate in the buffer across orbits, ultimately resulting in permanent data loss. The two-code scheme achieves an operational margin of $97.05\text{ kbit}$ ($2.11\%$ of pass capacity), providing sufficient headroom for operational survival.

\subsection{Unconsidered Effects, Operational Risks, and Mitigation Testing}

While the arithmetic identifies the two-code scheme as the sole viable option, several practical assumptions could compromise the link:

\begin{enumerate}
    \item \textbf{Mode Shift and Non-Stationary Distributions:} The analysis assumes exact $65\% / 35\%$ orbit partitioning with stationary event probabilities inside each mode. Seasonal beta-angle variations, extended eclipse duration, or spacecraft tumbling will alter these distributions. If eclipse duration increases, the overall bit rate shifts toward $L_{\mathrm{eclipse}} = 2.7200\text{ bits/record}$, pushing total bits past $4.608\text{ Mbit}$.
    \begin{itemize}
        \item \textit{Testing:} Perform Monte Carlo orbit simulations over varying beta angles and thermal scenarios to map worst-case telemetry volume against satellite orbital precession.
    \end{itemize}

    \item \textbf{Synchronization Loss and Bit Error Propagation:} Variable-length prefix codes are vulnerable to bit errors over noisy links (e.g., UHF fading). A single corrupted bit in a Huffman stream can cause loss of frame synchronization and catastrophic burst errors across subsequent records until the next frame boundary.
    \begin{itemize}
        \item \textit{Testing:} Inject bit-error profiles (e.g., Markov-channel models representing fading and noise bursts) into telemetry streams to evaluate frame error rates (FER) and test self-synchronization behavior.
    \end{itemize}

    \item \textbf{Buffer Overhead and Framing Synchronization:} The calculations omit physical-layer framing overheads such as sync markers (preambles), packet sequence numbers, and Forward Error Correction (FEC / CRC) parity bits.
    \begin{itemize}
        \item \textit{Testing:} Conduct end-to-end Hardware-in-the-Loop (HITL) RF testing between the CubeSat communications board and ground station hardware to measure net protocol overhead.
    \end{itemize}
\end{enumerate}

\end{document}